\section{连续统的观点}

我们想用连续统语言表述我们所做的一切，因为我们主张这一基本现象不是格点赝象。所以，我们需要用连续统语言理解涂抹意味着什么、涂抹后的可观测量是什么，以及如何用微扰与非微扰的连续统方法计算它们。接着，我们至少需要考虑大 $N$ 相变如何纳入连续统描述，以及在幺正场论中它们是否完全可以接受。

\subsection{涂抹的连续统威尔逊圈与第五维的暗示}

我们的第一个任务是理解涂抹在连续统中的含义。答案是：当噪声置零时，量 $f n$ 扮演朗之万（Langevin）方程中第五个``时间'' $\tau$ 的角色：
\begin{equation}
\frac{\partial A_\mu}{\partial\tau} =D_\mu F_{\mu\nu}
\label{langev}
\end{equation}
对小的 $f n$ 和小的 $\tau$，把方程 (\ref{eq:smearA})--(\ref{eq:fnsol}) 与 (\ref{langev}) 加以比较即可看清这一点。

若在右端加上一个 Zwanziger 项 \cite{zwa135}（它对我们的威尔逊圈或任何其他规范不变可观测量都没有影响）并把它移到左端，运动方程便取得五维形式，其中 $x_5\equiv
\tau$：

\be
F_{5\nu} =D_\mu F_{\mu\nu} \ee


当在右端加入噪声项时就得到规范理论的 Parisi-Wu 随机量子化 \cite{stoch13}。
这一框架已被证明对应于一个拓扑的、规范不变的五维理论 \cite{balzwa14}，其中第五个方向扮演特殊角色（例如它的单位是长度的平方，不同于其他四个坐标的单位）。欧氏转动不变性只在四维中存在。主体（bulk）理论在拓扑场论的意义上是平凡的，但边界上有非平凡动力学。这一框架令人想起全息的概念。
噪声项是这个构造的本质成分；我们的工作能否与 \cite{balzwa14} 的工作建立起更精确的关系，是一个开放的、也许有趣的问题。

在本文中我们只考察了处于一个共同固定 $\tau$ 的圈。
这个 $\tau$ 决定了圈的厚度，我们把它取为均匀的。但是，考虑 $\tau$ 绕圈变化的圈算符也完全有意义，它描述一根厚度变化的弦。而且，如果寻求一个绕圈处处局域的圈方程，就必须考虑这样的圈。的确，普通的四维运动方程进入了 $\tau$ 演化，因此有可能通过同时考虑每种可能厚度的圈，恢复我们在停留在四维中单纯把圈加粗时所丧失的 Migdal-Makeenko 方程的局域性。

反过来，Migdal-Makeenko 方程或许是通向基本的（无限细）大
$N$ QCD 弦的正确出发点；这根弦如今将生活在五个扭曲的维度中，并可能提供与四维肥弦所提供的描述等价的四维物理描述。这一思辨提示我们：我们那个要在四维中得到作为幺正矩阵的良定义连续统威尔逊圈的简单目标，把我们引向一个额外的维度并非纯属巧合。


\subsection{瞬子闭合能隙}

对小圈，绕它的平行输运应当接近恒等矩阵，在 $-1$ 附近没有本征值。
在无穷大 $N$，预期远离 $1$ 的本征值被完全压低，存在一个大的能隙。在矩阵模型中，对有限但大的
$N$，这种压低典型地随 $N$ 指数地小。

考虑平面内半径为 $R$ 的圆和一个中心放在圆周上某处的瞬子。
我们可以容易地计算该圆的威尔逊圈算符的本征值，并把圈在瞬子中心处剪开。

正规规范中的瞬子规范场为：
\be
A_\mu=\imath\frac{x^2}{x^2+\rho^2} U^\dagger \partial_\mu U,
\ee
其中
\be
U=\frac{x_4+\imath\vec x\cdot\vec\sigma}{\sqrt{x^2}}.
\ee

把我们的圆定义为 $(x_4-R)^2+x_3^2=R^2$，且 $x_{1,2}=0$，我们发现：
\be
{\cal P} e^{\imath\oint A_\mu dx_\mu}=e^{\imath\pi
(1-\frac{\rho}{\sqrt{\rho^2+4R^2}})\sigma_3 }  \ee



只要瞬子足够小且足够靠近圆周，它就会在
$-1$ 处产生一些本征值密度。对 $\Lambda R <<1$，大瞬子无关紧要，所以如果圈很小，就无需面对瞬子演算中著名的红外问题。这是一个 $e^{-N\cdot {\rm Const}}$ 效应，很可能与从任何矩阵模型普适（双重）标度函数得到的渐近行为光滑衔接。这也许足以挑选出矩阵模型微分方程的唯一解，而后者可以把我们一路带到大的 $R$，在那里有效弦理论与矩阵模型有一段交叠区域。显然，没有某些显式计算的话，这只是一个建议。

注意，正如我们定义的，涂抹的连续统版本不影响瞬子，因为它把~(\ref{langev}) 的右端置零。

\subsection{对大 $N$ 普适类的一个猜想}



对我们涂抹威尔逊圈中大 $N$ 相变的普适类，至少有两个明显的候选。如果只要求处理一个幺正矩阵在谱中张开能隙、且服从反映 CP 不变性的实条件的那种普遍模型，就会得到单个幺正矩阵模型和著名的 Gross-Witten 大 $N$
相变 \cite{gw12}。另一个直接的选项是定义在无量纲面积 $A$ 的球面上的二维 QCD，它具有 Douglas-Kazakov 相变。后者相变也发生在无限平面上的二维 QCD 中，不过这次不是对配分函数，而是对面积为 $A$ 的威尔逊圈算符；一种情形中球面上的相变与另一种情形中可观测量的相变在数学上是等同的。下面我们论证：无论在理论上还是数值上，第二个选项都更受青睐。

区分这两种情形可能有点令人困惑，因为两者都与二维 QCD 相关。Gross-Witten
相变与带单方块 Wilson 型作用量的格点 QCD 相关。
模型中的单个
幺正矩阵是绕一个基本方块的平行输运。作为格点耦合的函数，
这个矩阵在 $b$ 从零开始增大时于某临界值张开能隙，而这个矩阵的迹在相变点是非解析的。
现在考虑 Gross 和 Witten 的分析适用的格点模型，但考虑大于 $1\times 1$ 的圈。它的迹在 Gross-Witten 相变点仍非解析。但绕较大圈的平行输运将在大的 $b$ 值处张开谱隙，而在那个值处较大圈的迹并不表现出非解析行为。简言之，方块之迹的非解析行为是与连续统极限无关的格点效应。一旦我们同意避免看格点效应，那么显然我们应当宁可考虑把连续统二维 QCD 作为与我们为涂抹的四维威尔逊圈 ${\hat W}$ 找到的相变相关的大 $N$ 普适类的代表。

在连续统二维平面 QCD 中，对光滑简单圈，威尔逊圈算符的谱不依赖于
曲线的形状而只依赖于所围面积。当圈被标度时，唯一要紧的是面积变化；已知在无穷大 $N$ 其本征值分布在一个特定面积（以二维规范耦合常数为单位度量）处经历一次相变，恰好在本征值 -1 处张开能隙。然而，威尔逊圈算符任何有限次幂的迹当面积扫过谱的临界点时保持解析。二维中 $n$-缠绕威尔逊圈的期望值仍只依赖其基本面积——我们用 't Hooft 耦合把它无量纲化并记为 $A$。大 $N$ 的显式公式已知 \cite{KKwloop145}，为一个多项式与一个指数的乘积：
\begin{equation}
\frac{1}{N}\langle Tr { W^n } \rangle = \frac{1}{n} L_{n-1}^{(1)}
( 2A n) e^{-An}
\label{twodwil}
\end{equation}
对标度的依赖几乎不可能更解析了。


Gross-Witten 普适类与连续统 QCD 相关的那个普适类有联系但不必然等同。尽管如此，
两个候选都在双重标度极限中产生 Painleve II 方程 \cite{peris, gross-mat}。


另一个要记住的要点是：与 Gross-Witten 相变的普适类相联系的弦理论相当退化，看起来与任何设想的 QCD 弦性对偶显著不同 \cite{ias-double-cut}。
另一方面，二维 QCD 确实存在更具体的弦描述，源于 Gross 和 Taylor \cite{grosstay15}，它也有 sigma 模型表示 \cite{moorhorav155}。无限时空上连续统二维 QCD 威尔逊圈的大 $N$ 相变与定义在面积 $A$ 的球面上的平面二维 QCD 配分函数的 Douglas-Kazakov 相变属于同一类型 \cite{dougkaz157}。

还有一个有利于
二维 QCD 作为普适类的特征，它与瞬子有关。当我们取那个我们证明了的对四维小圈在 $-1$ 处本征值密度作出非零贡献的非阿贝尔瞬子解、而只看一个二维切片时，
就得到二维瞬子，文献~\cite{gross-mat} 中论证它们与 Douglas-Kazakov 相变有关。

所有这一切加在一起，使 Gross-Witten 相变作为我们相变普适类的代表处于不利地位，给我们留下二维平面 QCD 这一首选选项。威尔逊圈本征值分布的显式公式可在~\cite{duol} 找到，并且有可能~\cite{inst2d}
把这个公式与 (\ref{twodwil}) 联系起来。

\begin{figure}[ht]
\centering
\includegraphics[width=\textwidth]{images/durolefit.pdf}
\caption{四种不同威尔逊圈尺寸 $lt_c=0.740,0.660,0.560,0.503$ 的分布对 Durhuus-Olesen 分布的拟合。描述连续曲线的相关面积（Durhuus-Olesen 记号中的 $k$）分别由 $k=4.03,2.30,1.41,1.15$ 给出。 }
\label{durolefit}
\end{figure}


为了看到我们的猜想并未完全离谱，我们展示一幅对比图：把我们测得的几个本征值分布与 Durhuus 和 Olesen \cite{duol} 的精确无穷大 $N$
连续统解作比较：

图~\ref{durolefit} 显示了如何扫过临界尺寸；可以直接看到相变的效应，而且 Durhuus-Olesen 分布似乎在相变两侧都工作得相当好。尝试拟合 Gross-Witten 分布将会失败，因为在无隙一侧
Gross-Witten 分布总会在图~\ref{durolefit} 中
$\theta=\pm\frac{\pi}{2}$ 处穿过值
$\frac{1}{2}$。


支持该普适类的另一点从
equation~\ref{twodwil} 显而易见：对大面积，前置因子有振荡的符号。这些正是我们在图~\ref{fourier} 的傅里叶系数中看到的振荡符号。这些振荡符号可以被当作一个迹象：多次缠绕的圈的有效弦表示需要某种统计性质不同于玻色子的基本二维激发~\cite{gross-ooguri}。人们通常使用的简单有效弦模型并不处理多次缠绕的圈。


由于大 $N$ 普适类对我们纲领的重要性，我们现在转向用矩阵模型对其看似基本代表族的更一般的描述：

设想 $n$ 个 $SU(N)$ 矩阵，按照共轭不变的测度
$d\mu(U_i)$ 独立同分布。该测度在 $U_i=1$ 处有峰值且属单迹类型。作为每个本征值的函数，它沿圆周从 $1$ 走到 $-1$ 时单调下降，并在 $-1$ 处取最小值。人们定义的连续统威尔逊圈算符 $\hat W$，在其本征值分布于 $-1$ 张开能隙的大 $N$ 相变点附近，行为类似于 $W=\prod_i^n U_i$。测度的参数和 $n$ 可以调节到 $W$ 本征值分布在无穷大 $N$ 的各种临界行为，并很可能产生各种双重标度极限，对应于稳定性递减的大 $N$ 不动点。${\hat W}$ 的形状与尺度在相变附近解析地映入这些参数。最合理的是，最普遍、最稳定的相变才是相关的那个。

这个随机矩阵模型精确可解，我们现在概述其解法，
细节留待将来：

引入 $n$ 对 $N$ 分量 Grassmann 变量
$\bar\psi_i, \psi_i$，并通过简单递归证明
\be
\int \prod_i^n ( d\bar\psi_i d\psi_i ) e^{z\sum_i^n
\bar\psi_i\psi_i - \sum_i^n \bar\psi_i U_i\psi_{i+1}}=\det(z^n + W )
\ee
以上采用约定 $\psi_{n+1}\equiv - \psi_1$。

现在可以为 $SU(N)$ 不变双线性
$\bar\psi_i \psi_i$ 与 $\bar\psi_{i}\psi_{i+1}$ 引入辅助场。对 $U_i$ 变量的平均可以独立完成，剩下的积分在无穷大 $N$ 可用通常的鞍点方法处理。


在无限二维格点上选取轴向规范，容易看到：无论是在格点上还是在连续统中 \cite{duol}，二维平面 QCD 的非极小威尔逊圈确实由上述普适类描述。
类似地，Douglas-Kazakov 临界点也在这个普适类中，对应于最稳定的临界点。

直观上，这个普适类可以想象为以下阿贝尔情形的非阿贝尔推广：
给定光滑圈 $C$，令 $S$ 是以 $C$ 为边界的、拓扑为圆盘的曲面。在曲面 $S$
上定义一个二维 $U(1)$ 规范场，其通过铺砌 $S$ 的小面块的通量独立同分布。绕 $C$ 平行输运的相位便是所有这些通量之和。该和的分布一般受中心极限定理支配。显然，我们有了这一安排的非阿贝尔
推广~\cite{janik}，属于 Voiculescu \cite{voicu} 研究的类型。
